\section{Related work}

The closest theorem-level comparator is \citet[Assumptions 1--3 and Theorem 1(ii), p. 12]{Dorn2026GlobalOverlap}. In a fixed setup with uniform covariates on a cube, Gaussian outcomes, and Hölder response curves, satisfying Dorn's retained restrictions and local anti-concentration condition, Dorn establishes an essential-supremum probability upper bound uniformly over the source family. Writing \leanref{sym:d}{\ensuremath{d}} for the covariate dimension, \leanref{sym:beta}{\ensuremath{\beta}} for response smoothness, \leanref{sym:gamma}{\ensuremath{\gamma}} for the overlap exponent, and \leanref{sym:n}{\ensuremath{n}} for sample size, the pure-power rate in this paper's notation is
\[
n^{-\beta/(2\beta+d\gamma/(\gamma-1))}.
\]
For each probability tolerance, the source result supplies a finite threshold constant within that fixed setup. Dorn describes a count-based grid and bandwidth estimator computable from the known smoothness and observed treatment design without the overlap exponent \citep[Lemmas 11 and 13, pp. 23--24; proof, pp. 27--28]{Dorn2026GlobalOverlap}. The probability upper guarantee fixes the source family, its common constants, and the exponent before selecting the estimator and threshold constant. The construction is a direct antecedent for the present mesh selector; the new uniformity statement concerns one estimator and common risk constants across a prescribed compact exponent interval.

The present comparison concerns the guarantee and its model domain. Under the \leanref{S-1}{fixed model indices and class constants}, \cref{obj:thm:adaptive-minimax} gives an expected spatial-supremum guarantee over the envelopes in \cref{obj:def:model}, uniformly across a compact range of overlap exponents, with response smoothness fixed. These envelopes impose a bounded treated conditional mean and uniformly conditional sub-Gaussian residuals, intrinsic cube Hölder smoothness, a covariate density lower bound, and a global propensity-tail bound. Bounded potential outcomes are used for the testing subclass and the thin-strip example, rather than for the full upper-bound class. Dorn's local anti-concentration condition additionally requires a common positive fraction of the covariate mass in every sufficiently small neighborhood to have propensity comparable to that neighborhood's propensity supremum. Here, the global tail and the density lower bound support the equal-cell construction under the conditions in \cref{obj:ass:density,obj:ass:global-tail}. The distinction in admissible local treatment geometry is made precise in \cref{obj:def:dorn-a3,obj:prop:strict-enlargement}, while \cref{obj:prop:dorn-a3-free-corollary} states the corresponding Gaussian-envelope upper guarantee.

Degenerate-design regression supplies a complementary comparison at a single location. In a one-dimensional Gaussian regression experiment, \citet[Assumption M and Theorem 1, pp. 3--4]{Gaiffas2005DegenerateRates} derives pointwise minimax moment rates for a fixed design density that is locally symmetric and regularly varying around the evaluation point. The regression class is specified through a local polynomial approximation modulus, together with a bound on the function; regular variation of the design and modulus yields a power rate multiplied by a slowly varying factor. In the separate adaptive experiment, \citet[Abstract]{Gaiffas2005AdaptiveDegenerateDesign} constructs a local polynomial procedure adaptive to design degeneracy and Hölder smoothness, and identifies the logarithmic cost of that pointwise adaptation. These results characterize local information through the behavior of a specified design near one point. The present envelopes instead organize treatment-specific recovery through a global distributional restriction on the propensity.

Spatially normalized losses provide another way to express heterogeneous information. In one-dimensional Gaussian regression over Hölder classes, \citet[Assumption D, equation (2.1), and Theorems 1--2]{gaiffas2009} studies the supremum of estimation error divided by a location-dependent rate obtained from a local bias--variance balance. The adaptive upper result assumes a continuous design density that is positive everywhere or has finitely many zeros with specified local power behavior; the lower result uses an interval-mass condition and establishes optimality over intervals of prescribed shrinking size. Likewise, \citet[Sections 1 and 3]{schmidthieber2024} analyzes a weighted uniform criterion in one-dimensional Gaussian regression under a local doubling condition on the design distribution. Least squares over the unit-Lipschitz class attains the local rate for truths in a strictly smaller Lipschitz ball. Both comparisons retain the design's spatial variation in the loss normalization. Our criterion in \cref{obj:def:risks} measures expected unweighted spatial-supremum error, with worst-case difficulty indexed by the global overlap envelope.

The estimator combines familiar polynomial approximation with a particular treatment-design weighting rule. Local polynomial fitting approximates a smooth response within a neighborhood \citep{fan1996}; series methods approximate conditional expectations in finite-dimensional spaces, including polynomial and spline spaces \citep{newey1997}. Partitioning-based series estimation develops this approach through locally supported bases \citep{CattaneoFarrellFeng2020}. The construction in \cref{obj:def:template,obj:def:estimator} fits polynomials on partition cubes, assigns equal aggregate weight to the reference microcells, and selects its mesh through the treated-count feasibility rule in \cref{obj:def:mesh-selector}.

The causal estimand also determines the relevant comparison. \citet[Theorems 1--2 and Remark 11]{KennedyBalakrishnanRobinsWasserman2024} studies pointwise expected absolute error for the conditional average treatment effect under strict overlap, with separate Hölder restrictions on the treatment effect, propensity, and control response. Its attainable rates depend on their relative smoothness and on conditions controlling covariate-distribution estimation and the local fitting matrices. Our target is a treatment-specific response curve, with smoothness imposed directly on that response and information scarcity described by a global overlap tail. For functional estimation, \citet[Theorems 1--3]{MouDingWainwrightBartlett2023} derives local minimax mean-squared-error bounds for linear functionals over convex symmetric regression classes, together with attaining procedures for reproducing kernel Hilbert space classes under the stated noise, spectral, and sample-size conditions. That experiment measures accuracy through a specified functional and its observational coverage. The present minimax criterion evaluates recovery of the response throughout the covariate cube.
